% ======================================================================
\section{Related work}\label{sec:related}

\paragraph{Non-speech sound for DHH viewers.} Captioning practice~\cite{dcmp} describes a sound only when the picture
does not already show it; this is the visibility rule used here. Viewers want important sounds, their type and timing,
and a choice~\cite{may2025choices,alonzo2022beyond}; captions lose emotional information~\cite{revuelta2020eeg};
visual styling is welcome when it is readable and calm~\cite{delacerdapataca2024royale,prosodycaptions}. Sound-awareness
tools~\cite{jain2015hmd,homesound,soundwatch,protosound} show that users want urgent sounds fast and are bothered by
false alarms, which is why a wrong picture has a cost in the metric used here. None of this work generates pictures for ambient
sound in recorded video, chosen automatically by what is already visible.

\paragraph{Sound event detection.} Almost all sound taggers use the AudioSet label set~\cite{gemmeke2017audioset}
(527 classes in a tree, e.g.\ \emph{Bark} under \emph{Dog} under \emph{Animal}). Clip-level taggers such as
PANNs~\cite{kong2020panns}, AST~\cite{gong2021ast} and BEATs~\cite{chen2023beats} say \emph{what} is heard;
text-queried detectors such as CLAP~\cite{elizalde2023clap}, FLAM~\cite{wu2025flam} and FlexSED~\cite{hai2025flexsed}
score a sound name frame by frame. Detection is usually scored with event-based F1 and a time tolerance
(``collar'')~\cite{mesaros2016metrics,bilen2020psds}; the metric used here keeps the onset tolerance but reports a cost instead of
F1, in the spirit of cost curves~\cite{drummond2006costcurves}. A known weakness matters here: a quiet sound under speech
or music is ranked below \emph{Speech} or \emph{Music}.

\paragraph{Audio-language and vision-language models.} Audio-language models such as Qwen2-Audio~\cite{chu2024qwen2audio}
and Qwen3-Omni~\cite{xu2025qwen3omni} answer questions about audio; the system uses two of them as ``listeners''.
Vision-language models such as Qwen2.5-VL~\cite{bai2025qwen25vl} and LLaVA~\cite{liu2023llava} answer questions about
frames and are the only tool that can decide, for any sound, whether its source is visible. Their known failures matter:
a preference for some answer letters~\cite{zheng2024pride}, weak handling of negation~\cite{alhamoud2025negbench}, and
treating a visible object as the one that is sounding (Section~\ref{sec:analysis}).

\paragraph{Sound-source localisation and audio-to-image generation.} Localisation
models~\cite{senocak2023alignment,um2025objaware} find \emph{where} in the frame a sound comes from, but barely react to
off-screen or silent audio~\cite{juanola2025critical}, so they do not answer \emph{whether} it is visible.
Audio-to-image models~\cite{sungbin2023sound2scene,biner2024sonicdiffusion,ge2026audiocanvas} draw a picture from a whole
soundtrack, so they draw the scene rather than the one unseen source and cannot be told to leave a visible object out.
A text-to-image model (Qwen-Image~\cite{wu2025qwenimage}) with a short, guarded prompt is therefore used, and the
value of the on-screen check is measured against a \emph{blind} baseline that uses the same detector and generator but draws every
sound. Model judges~\cite{zheng2023judge} were tried for evaluation and failed (Appendix~\ref{app:judge}).
